\chapter{Vues JSP/JSTL et interface}\label{chap:vues}

\section{Principes}
\begin{itemize}
  \item \textbf{Aucun scriptlet.} Le \code{web.xml} déclare \code{scripting-invalid} : une balise \code{<\% \%>} dans une JSP fait
        échouer la compilation. Les pages n'utilisent que le langage d'expression (EL) et JSTL, ce qui empêche de glisser de la
        logique dans la vue.
  \item \textbf{Échappement systématique.} Toute valeur venue de l'utilisateur ou de la base passe par \code{<c:out>} ou par
        un gabarit qui l'utilise. C'est la défense principale contre l'injection de script (XSS).
  \item \textbf{Pas de dépendance externe.} jQuery est livré dans l'application, la police et les icônes sont locales
        (icônes en SVG intégré). L'application fonctionne sans réseau.
  \item \textbf{Fonctionne sans JavaScript.} Chaque action est un vrai formulaire POST ; le script ne fait qu'améliorer le
        parcours (panier sans rechargement).
\end{itemize}

\section{Gabarits (tag files)}
Les éléments répétés sont des fichiers de balises dans \code{WEB-INF/tags}, appelés avec le préfixe \code{t:}. Une modification
de mise en page se fait en un seul endroit (tableau~\ref{tab:tags}).

\begin{table}[htbp]
\centering
\caption{Gabarits JSP}\label{tab:tags}
\begin{tabularx}{\linewidth}{@{}lX@{}}
\toprule
\textbf{Balise} & \textbf{Rôle}\\
\midrule
\code{t:layout} & Page complète : barre latérale, entête mobile, message éphémère, tiroir du panier, pied de page.\\
\code{t:champ} & Champ de formulaire avec étiquette, aide et message d'erreur lié au champ.\\
\code{t:csrf} & Champ caché portant le jeton anti-CSRF.\\
\code{t:visuel} & Image du produit, ou dégradé coloré avec initiales quand il n'y en a pas.\\
\code{t:prix}, \code{t:stock} & Formatage du prix (« 129,90 MAD ») et étiquette de stock (disponible, dernières unités, épuisé).\\
\code{t:icone}, \code{t:vide} & Icônes SVG et page d'état vide avec illustration.\\
\bottomrule
\end{tabularx}
\end{table}

\section{Les pages}
Dix vues couvrent l'application : liste et fiche produit de la vitrine, connexion, inscription, panier, liste et formulaire
des produits, liste et formulaire des catégories, page d'erreur. La page d'erreur sert pour les codes 403, 404, 405 et 500 et
n'affiche jamais de trace technique.

\fig{10-vitrine}{La vitrine : filtre par catégorie, recherche et grille de produits}{fig:vitrine}
\fig{11-detail-produit}{Fiche produit avec étiquette de stock et choix de la quantité}{fig:detail}

\section{Formulaires et validation}
Quand un formulaire contient des erreurs, la page est réaffichée avec les valeurs saisies (sauf les mots de passe) et un
message sous chaque champ fautif. Le test se fait côté serveur : les attributs HTML (\code{required}, \code{maxlength})
n'améliorent que le confort, puisqu'ils peuvent être contournés.

\fig{12-inscription-erreurs}{Inscription refusée : messages sous les champs invalides}{fig:inscription}

\section{Mini-panier en jQuery}
Un clic sur « Ajouter au panier » envoie le formulaire en AJAX. Le serveur répond par l'état complet du panier ; le script
reconstruit le tiroir latéral, met à jour le badge du nombre d'articles et affiche une notification. Les boutons $+$ et $-$
et « Retirer » du tiroir appellent la même servlet.

Deux décisions de sécurité structurent le script.
\begin{itemize}
  \item Le texte venu du serveur est inséré avec \code{.text()}, jamais comme HTML : même un nom de produit piégé ne s'exécute pas.
  \item Le jeton CSRF est lu dans une balise \code{meta} de la page et envoyé en en-tête \code{X-CSRF-Token}.
\end{itemize}

\fig{13-ajout-panier-tiroir}{Ajout au panier sans rechargement : tiroir latéral, badge et notification}{fig:tiroir}
\fig{14-page-panier}{Page du panier : lignes, quantités, total et récapitulatif}{fig:panier}
\fig{15-stock-insuffisant}{Refus d'une quantité supérieure au stock disponible}{fig:stock}

\section{Administration}
Un administrateur voit dans la barre latérale une section dédiée. La liste des produits indique l'état (en vente ou retiré)
et propose trois actions : modifier, retirer ou remettre en vente, supprimer. Retirer un produit de la vente le conserve en base
et dans les paniers existants ; supprimer est refusé si le produit s'y trouve encore.

\fig{16-admin-produits}{Administration : liste des produits et actions}{fig:adminproduits}

\section{Adaptation aux petits écrans}
En dessous de 860 pixels de large, la barre latérale devient un menu coulissant ouvert par un bouton, la grille de produits
passe à une colonne et les lignes du panier se réorganisent. Les animations sont désactivées pour les personnes qui
l'ont demandé dans leur système (\code{prefers-reduced-motion}).

\fig{17-vue-mobile}{Affichage sur petit écran}{fig:mobile}
